\scanpage{54}
% heading
Antipodisme

% top, lines 2–4
\[
\begin{aligned}
&\deleted{x\longmapsto 2\unsure{o|c}\,\pi(x)-x}\\
\underline a&=-\mathrm{id}+\deleted{\unsure{2\pi\otimes o}}2\pi\otimes o=-\mathrm{id}+\pi\otimes o'\\
-\underline a&=\mathrm{id}-\pi\otimes o'
\end{aligned}
\]

% upper right
\margin{\uline{Attention} il se peut que \unsure{$c''$|$c'$} n'existe pas ``à dist. finie'', mais c' existe\ldots}

% upper, lines 5–7
\uline{NB} $\underline a$ défini sss $c'$ défini i.e. sss $\rho'$ inversible (c'est le cas si $f_V$ lisse). Notons que $\overbrace{\rho' o'=\unsure{\widetilde{o}|\widetilde{\sigma}|\widetilde{c}}}^{\text{si }n\ \unsure{impair|pair}}$ partt $\ne0$ donc $o'$ partt $\ne0$ donc $\mathcal O_S o'=\mathcal O_S\widetilde{o}$.

% slanted, upper left
% slanted
\margin{\uline{NB} Si $n$ pair ou $2$ inv, alors : $\underline a$ défini \uline{et} est une réflexion \uuline{str.} \uline{orth.} $\Longleftrightarrow\sigma'\rho'$ inv, i.e. $\sigma'$ et $\rho'$ inv.}

% upper middle, lines 7–10
$\langle o',\pi\rangle=2$, donc $-\underline a$ est une réflexion, sauf si en certains pts on a car. $=2$, \uline{et} \deleted{\unsure{alors}} \deleted{ce qui} rang impair
\margin{$\bar m$ \unsure{conclusion|conditions} pour \unsure{les|tts} points où 2 inv. Mais en car 2 pour $n$ imp. on a $o'=0$ d'où $\underline a=\mathrm{id}$.}

% upper middle, line 11
$\underline a$ commute aux $\tau_i$, il inverse $o$, $c'$ etc.

% upper middle, line 12
\[
\underline a(s_0)=\underline a(a_{-1})=\deleted{\unsure{2o-s_0}}o'-s_0=\widetilde{s}_0
\]

% middle, lines 13–15
\uline{NB} Dans $\unsure{E|\mathcal E}$, l'ens. des pts fixes de $\underline a$ est \deleted{\unsure{$\mathcal O_S o=\mathcal O_S$}} la droite engendrée par $o'$ si $2$ inv -- si $2\,1_S=0$ i.e. $S$ de car. 2, alors \deleted{\unsure{o'|c'est l'}} ens des pts fixes est $V$.

% middle, lines 16–18
Dans $\mathrm P(E)$, l'ens des pts fixes est $\mathrm P(V)\cup\{o'\}$ -- qui se réduit à $\mathrm P(V)$ en car. 2

On a $s_0\ne\widetilde{s}_0=\underline a\,s_0$ dans $\mathrm P(E)=X$.

% middle lower, line 19
Considérons le \deleted{droite} plan engendré par $s_0$ et $\widetilde{o}$ \deleted{\unsure{(ou $o'$)}}

% lower, line 20
\[
g(\lambda,\mu)=f(\lambda s_0+\mu\widetilde{o})=\deleted{\unsure{\tfrac12 f\lambda}}\lambda\mu\underbrace{\varphi(s_0,\widetilde{o})}_{-\zeta_0=-\sigma}+\mu^2 f(\widetilde{o})=-\sigma\lambda\mu-\sigma'\rho'\mu^2=-\mu(\sigma\lambda+\sigma'\rho'\mu)
\]

% lower, line 21
1) \uline{$n$ pair} : $\sigma=\sigma'$, $\rho=2\rho'$\qquad $f(\lambda s_0+\mu\widetilde{o})=-\mu\sigma(\lambda+\rho'\mu)$

% lower left, lines 22–29
\margin{\uline{NB} Si $\sigma$ inv. \deleted{\unsure{pour}} \inserted{pour} $f(\lambda s_0+\mu\widetilde{o})=0$, $\mu$ inv. sss $\lambda+\rho'\mu=0$ i.e. $\lambda=-\rho'\mu$ i.e. $(\lambda,\mu)\mathrel{\#}(-\rho',1)$ donc $\lambda s_0+\mu\widetilde{o}$ multiple de $-\rho's_0+\widetilde{o}$ ou encore (si $\rho'$ inv) de $-s_0+\frac{\widetilde{o}}{\rho'}=-s_0+o'=\widetilde{s}_0$}

% lower right, lines 22–30
\[
\begin{aligned}
&\text{c'est }=0\text{ sss }\sigma=0\\
&\text{c'est }\ne0\text{ en chaque fibre }\overset{\unsure{\text{sss}|\text{ssi}}}{\Longleftrightarrow}\sigma\ \text{inv. i.e. }\sigma'\ \text{inv.}\\
&\deleted{\text{dans ce cas,}}\inserted{\Longleftrightarrow}\quad Q_X\cap D_0\ \text{est un rev. étale de}\\
&\text{degré 2 de }S,\text{ dont }s_0\text{ est une section,}\\
&[\underline a_X s_0\text{ est la section complémentaire,}\\
&\hspace{3em}\text{quand }\underline a_X\text{ est définie}]\\
&\Longleftrightarrow Q_X\not\supset D_0\ \text{en tte fibre}\\
&\hspace{3em}(\text{c'est le cas si }f_E\ \text{lisse})\\
&\Longleftrightarrow\psi(\widetilde{o})\ne0\ \text{chaque fibre}\Longleftrightarrow(\widetilde{o})^\perp=V\ \text{chaque fibre}
\end{aligned}
\]

% foot, line 31
2) \uline{$n$ impair}\qquad $\sigma=2\sigma',\ \rho=\rho',$\qquad $f(\lambda s_0+\mu\widetilde{o})=-\mu\sigma'(2\lambda+\rho'\mu)$

% foot, left
% archive sheet number: 53

